\section{Discussion}
\label{sec:discussion}

\subsection{Why Authentication Proofs Are Insufficient}

Authentication, confidentiality, and secure channels provide complementary
foundations for agent communication. Authentication addresses who is
communicating; confidentiality concerns who can learn the protected
information. Neither establishes whether the receiving user's policy
permits that communication. Correctly identifying a participant and
protecting a credential in transit do not establish the participant's
entitlement to receive or use that credential.

The SAGA case study makes this distinction concrete without identifying an
authentication failure. The Stage~1 results retain the three queried
authentication correspondences and token secrecy, while the authorization
diagnostic exposes an obligation absent from that model. Separately, the
consequence model permits message acceptance under an assumed disagreement
between intended denial and implementation permission. These results have
different scopes; the Stage~1 guarantees do not automatically transfer to
all later models or to deployments.

As formulated in Chapter~3, accepted contact needs an applicable policy
decision linked to the same participants and credential. Prior issuance
establishes credential provenance, but not permission to issue.
Authentication binds decisions to identities; authorization determines
what those identities may do.

% Sources: paper/sections/03_problem_formulation.tex;
% paper/sections/07_evaluation.tex, Authorization coverage gap and consequence model.

\subsection{Authorization as a Security-Critical Layer}

The case study motivates treating policy evaluation code as a security-critical
component. The relevant chain is policy specification, decision
implementation, and enforcement point. The specification defines permission;
the evaluator implements that meaning; enforcement uses the resulting
decision to release material, issue a credential, or accept communication.
Reviewing only the last check leaves the correctness of its authorization
input unresolved.

The confirmed inconsistency thus raises a question beyond the local coding
error: how does the evaluator's contract justify subsequent protocol actions?
That relationship remains a verification obligation, not a completed
refinement result or evidence of defects in other matchers.

The gate controls further distinguish enforcement stages. In the model,
receiver denial prevents issuance and acceptance even when Provider release
remains possible; Provider denial blocks the earlier release boundary.
Assessing message acceptance alone therefore does not settle the obligation
governing material release. The protected action and its enforcement point
must be explicit. These controls do not establish a deployment repair.

% Sources: paper/PAPER_STORY.md, Sections 1, 6, and 8;
% proofs/CLAIM_EVIDENCE_MATRIX.md, Claims 2 and 4;
% paper/sections/03_problem_formulation.tex, intermediate permission conditions.

\subsection{Formal Verification Coverage}

Formal verification makes assumptions and obligations precise. Authentication
correspondences constrain the origins of protocol events, while secrecy
queries constrain adversarial knowledge. Their value is preserved when an
authorization obligation falls outside their scope: the missing obligation
does not invalidate the properties that were actually established.

The distinction is particularly important in the diagnostic model, where
\texttt{PolicyAllow} is never emitted. Its failed correspondence identifies
missing authorization support in the model; it does not independently detect
the Python defect. In the consequence model, authorization divergence is
supplied as scenario facts. ProVerif analyzes what follows under those facts,
rather than calculating a concrete matcher decision. The implementation
observation and the formal reachability witness must therefore remain
separate evidence layers.

For an authorization claim, verification scope must account for the relevant
decision layer, either through an explicit representation or through a
justified connection to independently analyzed implementation behavior.
Merely naming an event after a policy decision does not supply that connection.
Likewise, a true acceptance correspondence requires interpretation alongside
reachability: in the denial controls, it holds vacuously. The issue is the
alignment of model, query, and claimed behavior, rather than a limitation
inferred about ProVerif from behavior the model does not represent.

% Sources: proofs/CLAIM_EVIDENCE_MATRIX.md, Claims 1 and 3--4;
% paper/sections/06_formalization.tex, modeling boundary and queries;
% paper/sections/07_evaluation.tex, formal verification results.

\subsection{Generalization Beyond SAGA}

SAGA is a case study, not a representative sample of all agentic systems.
The observed gap is relevant to systems where policy decisions are delegated
to software components and then consumed at a separate enforcement point.
Agent frameworks, API gateways, capability systems, and policy engines offer
contexts in which this separation may warrant examination. This relevance
concerns the verification question, not evidence of a shared vulnerability.

Applying the question elsewhere requires identifying local policy semantics,
the evaluator, and the action relying on its decision, with evidence connecting
them under that system's assumptions. We have not verified these systems for
the same inconsistency or consequence. The transferable perspective is to
examine the policy-to-enforcement connection; prevalence and impact beyond
SAGA remain unmeasured.

% Sources: paper/PAPER_STORY.md, case-study positioning;
% proofs/CLAIM_EVIDENCE_MATRIX.md, Overall Research Argument;
% paper/sections/07_evaluation.tex, Limitations.

\subsection{Limitations}

\paragraph{Matcher-to-formal mapping.}
The matcher inconsistency is confirmed at implementation level, but the formal
model abstracts its result. The semantic boundary between Python execution
and scenario facts remains. Executable policy semantics integration would
need to relate concrete rules, specificity, budget interpretation, and
participant direction to the formal state and observations. The present
combination of evidence does not establish that refinement relation.

\paragraph{Deployment evidence.}
No production deployment exploit is demonstrated. The study does not report
a remote exploit, an attack on a deployed service, or an experiment involving
malicious infrastructure. The symbolic witness follows the prescribed peer
processes and does not establish arbitrary attacker control. It also does not
show completion of a downstream application or tool action. These limits
qualify the communication impact, not the confirmed matcher discrepancy.

\paragraph{Model abstraction.}
Private transport abstracts an authenticated, confidential channel rather
than verifying TLS. The Provider-to-receiver handoff is also abstract.
Although issuance and acceptance are modeled, the complete token lifecycle
is not: expiry, revocation, policy updates, and quota consumption are outside
the evaluated properties. Fresh tokens do not establish session isolation
or injective agreement. The models therefore do not simulate a complete
deployment or establish the correctness of these omitted behaviors.

\paragraph{Single case study.}
The evidence concerns SAGA and the explicit scenarios evaluated here. It
provides neither a survey of policy implementations nor a measure of how
often such inconsistencies occur. Further studies of additional agent
systems are needed to assess the broader applicability of the analysis,
with separate implementation evidence and modeling assumptions for each.

% Sources: paper/sections/03_problem_formulation.tex, static-policy scope;
% paper/sections/05_attack.tex, Evidence Boundary;
% paper/sections/06_formalization.tex, Formalization Boundary;
% paper/sections/07_evaluation.tex, Limitations.
